逐位置使用同一个函数读不到上下文，把整个序列压成一个固定摘要又无法让不同的查询寻找不同的信息。注意力给出的第三条路是：保留全部可读位置，让每个查询根据当前内容临时算出自己的加权读取规则。连接结构不再由架构写死，而由内容当场决定。

这是一笔交易，收益明确，账单也明确。固定的连接结构原本免费提供了三样东西——它知道谁在前谁在后，它的深度是可控的，它的开销随长度线性增长。把它换掉之后，这三样都得另外补回来。本章的组织就是这笔交易的完整清单。

第一篇给出收益本身：Query、Key、Value 的分工，softmax 沿哪条轴归一化，mask 在声明什么。第二篇说明这份收益在同一笔参数预算里如何切分——多头不增加参数，它把匹配能力拆成若干条低秩通道，秩上界就是分工的价钱。第三到第五篇是三笔必须补交的账：注意力对输入置换严格等变，所以位置只能外挂；恒等通路一旦被归一化打断，深度就堆不上去；全局可读的标价是训练侧 \(n^2\) 的注意力矩阵与推理侧线性增长的 KV cache。第六篇退后一步，把注意力与它取代的循环路线放在同一组坐标上比较，说明两者在梯度路径、训练并行与推理成本上的优劣正好互补。

\subsection{怎么读这一章}

第一篇是主干，其余五篇都建立在它的记号上。若时间有限，读完第一篇后可以直接跳到最贴近当前困惑的一篇：调不动深模型看第四篇，显存爆了看第五篇，长度外推出问题看第三篇。

一个自检标准：能否说清每一行 softmax 在哪些位置上归一化、哪些位置被允许读取，以及去掉位置编码之后模型究竟丢失了什么。三个问题都能答上来，公式记不记得住并不重要。

\subsection{本章篇目}

以下篇目按依赖顺序排列，也可以从具体问题进入，再沿篇内链接回补。

\begin{itemize}
\tightlist
\item \hyperref[sec:14-Deep-Learning/05-Attention-and-Sequence-Models/01]{``Self-Attention 是内容决定的加权读取''}：把``读谁''从固定连接变成内容依赖的计算。
\item \hyperref[sec:14-Deep-Learning/05-Attention-and-Sequence-Models/02]{``多头不是冗余而是分工''}：单头输出是凸组合，一个头只能表达一种对齐；参数量不变，换来的是秩与通道数的权衡。
\item \hyperref[sec:14-Deep-Learning/05-Attention-and-Sequence-Models/03]{``位置必须外挂，因为注意力看不见顺序''}：置换等变的证明，以及正弦编码、可学习编码与 RoPE 各自打破对称的方式。
\item \hyperref[sec:14-Deep-Learning/05-Attention-and-Sequence-Models/04]{``残差与归一化让深层可训''}：恒等项买下的那条梯度通路，LayerNorm 为何不能换成 BatchNorm，Pre-LN 与 Post-LN 的真实权衡。
\item \hyperref[sec:14-Deep-Learning/05-Attention-and-Sequence-Models/05]{``二次复杂度与推理时的显存账''}：训练撞的是 \(n^2\)，推理撞的是 KV cache；稀疏、线性注意力、GQA 与 FlashAttention 各自卖掉了什么。
\item \hyperref[sec:14-Deep-Learning/05-Attention-and-Sequence-Models/06]{``序列建模的两条路线各自付什么代价''}：梯度路径、训练并行与推理成本三行对照，以及状态空间模型想同时拿到什么。
\end{itemize}
